% TOP_PROBLEM: {"id":30007668,"problem_number":"LOCAL-30007668","title":"When to decentralize government","category_id":21,"category":{"id":21,"name":"economics","display_name":"Economics","description":"Open economic theory statements, published research agendas, and proposed scoped specifications; question provenance and historical priority are qualified per record.","slug":"economics","order_index":21,"created_at":"2026-10-08T00:00:00Z"},"status":"research_question","external_url":null,"legacy_ids":[],"published":false,"statement":"In the explicitly specified setting: Under which information and capacity conditions does decentralization improve public services, and when does it magnify local capture, inequality, or coordination failures? The mathematical statement above fixes the target, assumptions and scope of this catalogue formulation.","economics":{"original_id":157,"field":"Political economy and institutions","kind":"design","formalization_basis":"adapted_research_question","source_status":"research_agenda","priority_status":"earliest_found","priority_note":"Earliest directly inspected qualifying passage in this audit; earlier origins were not established. A review or research-agenda statement does not imply original priority.","earliest_verified_reference":{"authors":["Timothy Besley","Robin Burgess","Adnan Khan","Guo Xu"],"title":"Bureaucracy and Development","year":2022,"url":"https://www.robinburgess.com/s/Bureaucracy_and_Development.pdf","doi":"10.1146/annurev-economics-080521-011950","locator":"Section 4.1, printed p. 413, paragraph beginning \"Although a hierarchical organization\"","evidence_summary":"The tradeoff between local adaptation and fragmented coordination is explicitly identified as an important research area."},"current_reference":{"authors":["Timothy Besley","Robin Burgess","Adnan Khan","Guo Xu"],"title":"Bureaucracy and Development","year":2022,"url":"https://www.robinburgess.com/s/Bureaucracy_and_Development.pdf","doi":"10.1146/annurev-economics-080521-011950","locator":"Section 4.1, printed p. 413, paragraph beginning \"Although a hierarchical organization\"","evidence_summary":"The tradeoff between local adaptation and fragmented coordination is explicitly identified as an important research area."},"formalization_note":"The cited passage establishes a research direction, not this exact mathematical statement. The notation, population, policy menu, assumptions, and estimands are an original adaptation for this catalogue; no universal unsolved theorem or source priority is claimed."},"statusQualification":"Published question or research agenda verified at the cited locator. The mathematical specification is adapted for this catalogue and includes additional assumptions; source publication does not establish first-ever priority or certify this exact model remains unsolved. The cited passage establishes a research direction, not this exact mathematical statement. The notation, population, policy menu, assumptions, and estimands are an original adaptation for this catalogue; no universal unsolved theorem or source priority is claimed.","statusReviewedAt":"2026-10-08"}
% FIRST_POSTED: 2026-10-08
% ENTRY_KIND: statement_only
% !TeX program = lualatex
\documentclass[11pt]{article}
\usepackage{fontspec}
\usepackage[a4paper,margin=25mm]{geometry}
\usepackage{amsmath,amssymb,mathtools}
\usepackage{xurl}
\usepackage[hidelinks]{hyperref}
\newcommand{\sref}[2]{\hyperlink{source:#1:#2}{[S#2]}}
\setlength{\emergencystretch}{3em}
\begin{document}
% problemId:30007668
\section*{When to decentralize government}\label{30007668}
\subsection{Definitions and mathematical statement}
Consider regions \(j\) within a specified state providing a public service. Governance rule \(d\in[0,1]\) specifies the degree of local authority over budgets, personnel, and service design, separately from the source of financing. Let \(Q_j(d)\) be service quality, \(K_j\) prereform local administrative capacity, \(I_j\) local information, and \(E_j\) prereform elite influence. Let \(C(d)\) be total cost and \(V(d)=\sum_j\omega_jQ_j(d)\) a declared population-weighted objective. Estimate heterogeneous policy effects \(E[Q_j(d)-Q_j(0)\mid K_j,I_j,E_j]\) and characterize the feasible rule maximizing \(V\) subject to a budget and explicit inequality constraints. Decentralization can improve local adaptation while reducing coordination, so include cross-region spillovers and national distributional effects. Information technology may change the value of local information and must be incorporated into the policy environment. Use credible reform or assignment variation and avoid treating selected high-capacity localities as representative. The cited review expressly identifies this organizational tradeoff as a research area. The proposed causal contrasts and constrained authority rule make the broad question precise without assuming decentralization has the same effect for every service or state.

\par\textbf{Formulation and scope.} The cited passage establishes a research direction, not this exact mathematical statement. The notation, population, policy menu, assumptions, and estimands are an original adaptation for this catalogue; no universal unsolved theorem or source priority is claimed.

\par\textbf{Open-question provenance.} An explicit question or research agenda was verified at the source locator. This does not by itself certify that every later version remains unresolved \sref{30007668}{1}.

\par\textbf{Historical priority.} Earliest directly inspected qualifying passage in this audit; earlier origins were not established. A review or research-agenda statement does not imply original priority.

\subsection{Short English statement}
In the explicitly specified setting: Under which information and capacity conditions does decentralization improve public services, and when does it magnify local capture, inequality, or coordination failures? The mathematical statement above fixes the target, assumptions and scope of this catalogue formulation.

\subsection{Sources}
\begin{itemize}\raggedright
\item[\textnormal{[S1]}]\hypertarget{source:30007668:1}{} Timothy Besley, Robin Burgess, Adnan Khan, Guo Xu. 2022. Bureaucracy and Development. Section 4.1, printed p. 413, paragraph beginning "Although a hierarchical organization".
\url{https://www.robinburgess.com/s/Bureaucracy_and_Development.pdf}.
DOI: 10.1146/annurev-economics-080521-011950.
{\small\textit{Earliest verified question/agenda reference; historical priority qualified below. The tradeoff between local adaptation and fragmented coordination is explicitly identified as an important research area.}}
\end{itemize}
\end{document}
